\begin{question}{Efficiently Simulating Bounded Stacks with RNNs}

\begin{subquestion}
    The language $D(2, 3)$ consists of brackets of $k = 2$ types with maximum nesting depth $m = 3$. 
    We encode the unclosed opening brackets using a stack whose \emph{rightmost} symbol is the top of the stack.
    $[\openBr{1}\openBr{2}]$ transitions to $[\openBr{1}]$ after reading $\closeBr{2}$.
    \begin{itemize}
        \item Reading a closing bracket $\closeBr{i}$: pop the top of the stack if it equals $\openBr{i}$; otherwise (mismatch or empty stack), reject.
        \item Reading an opening bracket $\openBr{i}$: push it on top, unless the stack already has $m = 3$ symbols, in which case the depth bound is exceeded and we reject.
    \end{itemize}
    \begin{table}[H]
    \centering
        \begin{tabular}{cccc}
            \toprule
            & $\stackseq_1 = \openBr{2}$ &
            $\stackseq_2 = \openBr{1} \openBr{1} \openBr{1}$ &
            $\stackseq_3 = \openBr{2} \openBr{1} \openBr{2}$ \\
            \midrule
            $\closeBr{2}$ & $\emptystack$ & reject & $\openBr{2}\openBr{1}$ \\
            % \midrule
            $\closeBr{1}$ & reject & $\openBr{1}\openBr{1}$ & reject \\
            % \midrule
            $\openBr{2}$ & $\openBr{2}\openBr{2}$ & reject & reject \\
            \bottomrule
        \end{tabular}
        \caption{The first row of the table represents different stack configurations, while the first column lists the newly read symbols.}
        \label{tab:stackconfigs}
    \end{table}
\end{subquestion}

\begin{subquestion}
    We construct a DFA $\mathcal{A} = \fsatuple$ that recognizes $D(k, m)$ by encoding every legal stack configuration as a state. 
    Let $\openBr{} = \{\openBr{1}, \dots, \openBr{k}\}$ and $\closeBr{} = \{\closeBr{1}, \dots, \closeBr{k}\}$.
    
    \textbf{Alphabet:} $\alphabet = \openBr{} \cup \closeBr{}$, so $\nsymbols = 2k$.
    
    \textbf{States:} $\states = \{[\stackseq] \mid \stackseq \in \openBr{}^{\le m}\} \cup \{\rejectState\}$, where $[\stackseq]$ corresponds to the stack of unclosed brackets $\stackseq$ with $|\stackseq| \le m$, and $\rejectState$ is an absorbing rejecting state. $\acceptState \defeq [\emptystack]$ for the empty-stack state.
    
    \textbf{Initial states:} $\initial = \{\acceptState\}$ (start with an empty stack).
    
    \textbf{Final states:} $\final = \{\acceptState\}$ (only the empty stack accepts).
    
    \textbf{Transition $\trans : \states \times \alphabet \to \states$:} For valid prefix $\stackseq \in \openBr{}^{*}$ with $[\stackseq] \in \states$ and $i, j \in \{1, \dots, k\}$:
    
    \begin{enumerate}
        \item \emph{Valid push} (if $|\stackseq| < m$): $\trans([\stackseq], \openBr{i}) = [\stackseq\, \openBr{i}]$.
        \item \emph{Valid pop}: $\trans([\stackseq\, \openBr{i}], \closeBr{i}) = [\stackseq]$.
        \item \emph{Depth overflow} (if $|\stackseq| = m$): $\trans([\stackseq], \openBr{i}) = \rejectState$.
        \item \emph{Mismatched close} (if $i \ne j$): $\trans([\stackseq\, \openBr{j}], \closeBr{i}) = \rejectState$.
        \item \emph{Underflow}: $\trans(\acceptState, \closeBr{i}) = \rejectState$.
        \item \emph{Sink}: $\trans(\rejectState, s) = \rejectState$ for every $s \in \alphabet$.
    \end{enumerate}
    These cases specify $\trans$ on every $(\states \setminus \{\rejectState\}) \times \alphabet$ pair as well as on $\{\rejectState\} \times \alphabet$, so $\trans$ is deterministic.

    \paragraph{Correctness.} A string $\str$ is in $D(k, m)$ iff (i) the brackets are balanced and (ii) the running depth $d(\str_{\le t}) \le m$.
    The state $[\stackseq_t]$ reached after reading $\str_{\le t}$ stores the multiset \emph{and order} of currently unclosed brackets, so $|\stackseq_t| = d(\str_{\le t})$. The transitions enforce the two conditions:
    \begin{itemize}
        \item Any depth violation triggers rule (3), sending $\mathcal{A}$ into the absorbing $\rejectState$.
        \item Any mismatched close (rule 4) or underflow (rule 5) likewise sends $\mathcal{A}$ to $\rejectState$.
    \end{itemize}
    Conversely, while $\str_{\le t} \in D(k, m)$ at every prefix, only rules (1) and (2) fire and the automaton stays inside $\{[\stackseq] \mid \stackseq \in \openBr{}^{\le m}\}$. Since $\final = \{\acceptState\}$, the run accepts iff after reading the whole string the stack is empty, i.e.\ all opens were correctly matched. Hence $\mathcal{L}(\mathcal{A}) = D(k, m)$.

    \paragraph{Complexity.}
    \[
        |\states| \;=\; 1 + \sum_{j=0}^{m} k^{j} \;=\; 1 + \frac{k^{m+1} - 1}{k - 1} \;=\; \mathcal{O}(k^{m}).
    \]
    Because $\mathcal{A}$ is a DFA, the number of transitions is
    \[
        |\trans| \;=\; |\states| \cdot |\alphabet| \;=\; \mathcal{O}(k^{m}) \cdot 2k \;=\; \mathcal{O}(k^{m+1}).
    \]
\end{subquestion}

\begin{subquestion}
    For $D(2, 3)$ we have $k = 2$ and $m = 3$, so each hidden state lives in $\mathbb{R}^{k \cdot m} = \mathbb{R}^{6}$.
    Writing $\mathbf{e}_1 = (1, 0)^\top$ and $\mathbf{e}_2 = (0, 1)^\top$ for the canonical basis vectors of $\mathbb{R}^{2}$ associated with $\openBr{1}$ and $\openBr{2}$, Eq.~(27) places the top of the stack in the first $k$-block and pads unused bottom blocks with zeros.
    The empty stack maps to $h([\emptystack]) = \mathbf{0}_{6}$, and the initial embeddings of $\stackseq_1, \stackseq_2, \stackseq_3$ are
    \begin{align*}
        h(\stackseq_1) &= h([\openBr{2}]) = \begin{pmatrix} \mathbf{e}_2 \\ \mathbf{0}_2 \\ \mathbf{0}_2 \end{pmatrix} = (0, 1, 0, 0, 0, 0)^\top, \\
        h(\stackseq_2) &= h([\openBr{1}\openBr{1}\openBr{1}]) = \begin{pmatrix} \mathbf{e}_1 \\ \mathbf{e}_1 \\ \mathbf{e}_1 \end{pmatrix} = (1, 0, 1, 0, 1, 0)^\top, \\
        h(\stackseq_3) &= h([\openBr{2}\openBr{1}\openBr{2}]) = \begin{pmatrix} \mathbf{e}_2 \\ \mathbf{e}_1 \\ \mathbf{e}_2 \end{pmatrix} = (0, 1, 1, 0, 0, 1)^\top.
    \end{align*}
    Re-encoding the resulting stacks from part~(a) via Eq.~(27) yields the embeddings:
    \begin{table}[H]
    \centering
    \small
        \begin{tabular}{cccc}
            \toprule
            & $\stackseq_1 = \openBr{2}$ &
            $\stackseq_2 = \openBr{1} \openBr{1} \openBr{1}$ &
            $\stackseq_3 = \openBr{2} \openBr{1} \openBr{2}$ \\
            \midrule
            $\closeBr{2}$ & $(0, 0, 0, 0, 0, 0)^\top$ & reject & $(1, 0, 0, 1, 0, 0)^\top$ \\
            % \midrule
            $\closeBr{1}$ & reject & $(1, 0, 1, 0, 0, 0)^\top$ & reject \\
            % \midrule
            $\openBr{2}$ & $(0, 1, 0, 1, 0, 0)^\top$ & reject & reject \\
            \bottomrule
        \end{tabular}
        \caption{The first row of the table represents different stack configurations, while the first column lists the newly read symbols.}
        \label{tab:stackembeddings}
    \end{table}
\end{subquestion}

\begin{subquestion}
    Let $\mathbf{I}_k$ and $\mathbf{0}_k$ denote the $k\times k$ identity and zero matrices, and write $\mathbf{h}$ as $m$ stacked $k$-blocks where block~$j$ is the $j$-th symbol from the top of the stack. 
    Block~1 is the top, blocks beyond the stack depth are zero.
    Pushing onto the stack means \emph{shifting all $k$-blocks one position down} and writing the new symbol into block~1. 
    Popping means \emph{shifting all $k$-blocks one position up} and zeroing block~$m$. 
    Realize this with two block-banded $km \times km$ matrices,
    \[
        \mathbf{U}_{\text{push}} \;=\;
        \begin{pmatrix}
            \mathbf{0}_k & \mathbf{0}_k & \cdots & \mathbf{0}_k & \mathbf{0}_k \\
            \mathbf{I}_k & \mathbf{0}_k & \cdots & \mathbf{0}_k & \mathbf{0}_k \\
            \mathbf{0}_k & \mathbf{I}_k & \cdots & \mathbf{0}_k & \mathbf{0}_k \\
            \vdots & \vdots & \ddots & \vdots & \vdots \\
            \mathbf{0}_k & \mathbf{0}_k & \cdots & \mathbf{I}_k & \mathbf{0}_k
        \end{pmatrix},
        \qquad
        \mathbf{U}_{\text{pop}} \;=\;
        \begin{pmatrix}
            \mathbf{0}_k & \mathbf{I}_k & \mathbf{0}_k & \cdots & \mathbf{0}_k \\
            \mathbf{0}_k & \mathbf{0}_k & \mathbf{I}_k & \cdots & \mathbf{0}_k \\
            \vdots & \vdots & \vdots & \ddots & \vdots \\
            \mathbf{0}_k & \mathbf{0}_k & \mathbf{0}_k & \cdots & \mathbf{I}_k \\
            \mathbf{0}_k & \mathbf{0}_k & \mathbf{0}_k & \cdots & \mathbf{0}_k
        \end{pmatrix},
    \]
    and define the recurrence-matrix function $\mathbf{U} : \overline{\alphabet} \to \mathbb{R}^{km \times km}$ by
    \[
        \mathbf{U}(y_t) \;=\;
        \begin{cases}
            \mathbf{U}_{\text{push}} & \text{if } y_t \in \{\openBr{1}, \dots, \openBr{k}\}, \\
            \mathbf{U}_{\text{pop}}  & \text{if } y_t \in \{\closeBr{1}, \dots, \closeBr{k}\}, \\
            \mathbf{I}_{km}          & \text{if } y_t = \eos.
        \end{cases}
    \]
    The choice $\mathbf{U}(\eos) = \mathbf{I}_{km}$ keeps the stack representation intact at termination, since $\eos$ is not a stack operation. 
    The input matrix $\mathbf{V} \in \mathbb{R}^{km \times (2k+1)}$ injects the new opening bracket into block~1 only when $y_t$ is an opening bracket, so we set
    \[
        \mathbf{V} \;=\;
        \begin{pmatrix}
            \mathbf{I}_k \\
            \mathbf{0}_k \\
            \vdots \\
            \mathbf{0}_k
        \end{pmatrix}
        \begin{pmatrix} \mathbf{I}_k & \mathbf{0}_k & \mathbf{0}_{k\times 1} \end{pmatrix}
        \;=\;
        \begin{pmatrix}
            \mathbf{I}_k & \mathbf{0}_k & \mathbf{0}_{k\times 1} \\
            \mathbf{0}_k & \mathbf{0}_k & \mathbf{0}_{k\times 1} \\
            \vdots & \vdots & \vdots \\
            \mathbf{0}_k & \mathbf{0}_k & \mathbf{0}_{k\times 1}
        \end{pmatrix},
    \]
    where the columns are ordered as $\openBr{1}, \dots, \openBr{k}, \closeBr{1}, \dots, \closeBr{k}, \eos$. 
    With this $\mathbf{V}$, $\mathbf{V}\mathbf{o}_{\openBr{i}} = (\mathbf{e}_i, \mathbf{0}_{(m-1)k})^\top$ and $\mathbf{V}\mathbf{o}_{\closeBr{i}} = \mathbf{V}\mathbf{o}_{\eos} = \mathbf{0}_{km}$.

    \paragraph{Correspondence with $\mathcal{A}$.} Suppose $\mathbf{h}_{t-1}$ encodes a non-rejecting state $q = [\openBr{i_1}\cdots\openBr{i_{m'}}] \in \states$ via Eq.~(27). We check the three transition kinds.

    \emph{Valid push.} If $y_t = \openBr{i}$ and $m' < m$, $\mathcal{A}$ moves to $q' = [\openBr{i_1}\cdots\openBr{i_{m'}}\openBr{i}]$. 
    Multiplying out, $\mathbf{U}_{\text{push}}\mathbf{h}_{t-1}$ takes block~$j$ of $\mathbf{h}_{t-1}$ into block~$j{+}1$ for $j = 1, \dots, m{-}1$ and zeros block~1. 
    Adding $\mathbf{V}\mathbf{o}_{\openBr{i}}$ writes $\mathbf{e}_i$ into block~1. 
    The resulting blocks are $(\mathbf{e}_i, \mathbf{e}_{i_{m'}}, \dots, \mathbf{e}_{i_1}, \mathbf{0}, \dots, \mathbf{0})$, exactly $h(q')$.

    \emph{Valid pop.} If $y_t = \closeBr{i}$ with $i = i_{m'}$, $\mathcal{A}$ moves to $q' = [\openBr{i_1}\cdots\openBr{i_{m'-1}}]$. 
    Now $\mathbf{U}_{\text{pop}}\mathbf{h}_{t-1}$ takes block~$j$ into block~$j{-}1$ for $j = 2, \dots, m$ and zeros block~$m$, while $\mathbf{V}\mathbf{o}_{\closeBr{i}} = \mathbf{0}$. 
    The resulting blocks are $(\mathbf{e}_{i_{m'-1}}, \dots, \mathbf{e}_{i_1}, \mathbf{0}, \dots, \mathbf{0})$, exactly $h(q')$.

    \emph{$\eos$.} The DFA does not consume $\eos$ (cf.\ part (b)).
    We simply preserve the stack via $\mathbf{U}(\eos) = \mathbf{I}_{km}$ and $\mathbf{V}\mathbf{o}_{\eos} = \mathbf{0}$, so $\mathbf{h}_t = \mathbf{h}_{t-1}$.

    \emph{Invalid transitions.} Pushing on a full stack ($m' = m$) loses block~$m$ (shifted out), and popping a mismatched or empty stack discards the wrong $\mathbf{e}$ or stays at $\mathbf{0}$. 
    In either case $\mathbf{h}_t$ no longer encodes valid stack. The dynamics map alone cannot detect this. 
    It is the recognition module of part~(f) that assigns probability~$<\eta$ to such continuations, so in the locally $\eta$-truncated support these strings are rejected.

    Hence, on every \emph{valid} prefix, the RNN's hidden states are in one-to-one correspondence with the states $\mathcal{A}$ visits, and the recurrence $\mathbf{h}_t = \mathbf{U}(y_t)\mathbf{h}_{t-1} + \mathbf{V}\mathbf{o}_{y_t}$ implements $\trans$ exactly.
\end{subquestion}

\begin{subquestion}

\end{subquestion}

\begin{subquestion}

\end{subquestion}

\begin{subquestion}

\end{subquestion}

\begin{subquestion}

\end{subquestion}

\begin{subquestion}

\end{subquestion}

\end{question}